\section{多任务强化学习}

\subsection{动机：为什么 RL 迟早要走向 generalist policy}

讲者把多任务 RL 的动机说得非常直白：真实世界中的智能系统几乎从不只面对单一任务。通用助手要处理旅行预订、购物、问答；机器人要走路、下蹲、搬运、清理；推荐系统要面对不同用户偏好。真正的问题不是 ``某个策略能否学会一个任务''，而是 ``我们能否训练出一套能够在许多任务上共享知识的策略''。

\begin{figure}[H]
\centering
\includegraphics[width=0.9\textwidth]{slides-images/slide-11.jpg}
\caption{多任务 RL 的应用图景：从 LLM 助手到移动操作机器人与推荐系统}
\end{figure}

\begin{figure}[H]
\centering
\includegraphics[width=0.9\textwidth]{slides-images/slide-12.jpg}
\caption{多任务 RL 的深层动机：把数据复杂度在任务之间摊薄，构建更 general 的系统}
\end{figure}

\subsection{任务其实就是一组不同的 MDP}

\begin{importantbox}{多任务 MDP 的形式化}
将多任务问题写成一组 MDP $\{\mathcal{M}_i\}$，每个任务都可能在状态空间、动作空间、初始状态分布、动力学和奖励函数上发生变化。实际训练时，常把任务标识 $z_i$ 并入状态，写成：
$$\tilde{\mathcal{M}} = (\mathcal{S} \times \mathcal{Z}, \mathcal{A}, \tilde{p}, \tilde{r}, \gamma)$$
于是策略就变成条件策略 $\pi(a\mid s, z)$。
\end{importantbox}

这个定义比很多直觉化的 ``任务标签'' 更宽泛。一个任务不一定只是 ``做不同事情''，也可能是不同用户、不同初始配置、不同动力学设置，甚至不同身体结构。

\begin{figure}[H]
\centering
\includegraphics[width=0.9\textwidth]{slides-images/slide-14.jpg}
\caption{任务分布的例子：推荐系统、角色动画、多机器人 RL 的任务差异}
\end{figure}

\subsection{任务标识如何提供给策略}

任务描述 $z$ 可以采用多种形式：
\begin{itemize}
    \item \textbf{One-hot task ID}：实现简单，但几乎不支持零样本泛化。
    \item \textbf{自然语言指令}：更适合开放词汇任务，也更容易与 VLA/LLM 结合。
    \item \textbf{目标状态}：把任务写成 ``去到哪里''。
    \item \textbf{奖励参数或视频示范}：适合任务有更复杂语义时使用。
\end{itemize}

\begin{figure}[H]
\centering
\includegraphics[width=0.9\textwidth]{slides-images/slide-15.jpg}
\caption{任务标识可以是索引、语言、视频示范或目标状态}
\end{figure}

\begin{figure}[H]
\centering
\includegraphics[width=0.9\textwidth]{slides-images/slide-16.jpg}
\caption{把任务标识并入状态后，多任务 RL 可被重写为标准 MDP}
\end{figure}

\begin{knowledgebox}{这个重写的重要性}
一旦接受 $s=(\bar{s}, z)$ 的写法，多任务 RL 就不再需要一整套新的 RL 数学。单任务里的 policy、Q-function、critic 基本都可以直接条件化改写：
\begin{itemize}
    \item $\pi_\theta(a\mid s)\rightarrow \pi_\theta(a\mid \bar{s}, z)$
    \item $Q_\phi(s,a)\rightarrow Q_\phi(\bar{s}, a, z)$
\end{itemize}
这也是多任务 RL 在工程实现上远没有表面那么陌生的原因。
\end{knowledgebox}

\subsection{多任务 imitation learning 的具体做法}

最直接的多任务学习方式仍然是行为克隆：收集多个任务的演示数据，训练条件策略 $\pi_\theta(a\mid s,z)$。但课程特别提醒了一个经常被忽视的实践技巧：\textbf{stratified sampling}。也就是让每个 mini-batch 尽量包含多个任务的数据，避免训练被高频任务主导。

\begin{figure}[H]
\centering
\includegraphics[width=0.9\textwidth]{slides-images/slide-18.jpg}
\caption{多任务 imitation learning：条件策略配合分层采样，降低梯度方差}
\end{figure}

\begin{knowledgebox}{多任务学习为什么能共享数据}
不同任务之间经常共享技能子结构。例如 ``拿起红方块'' 和 ``拿起蓝方块'' 共享抓取动作；``擦桌子'' 和 ``收拾厨房'' 共享移动与定位能力。多任务学习的本质价值，就是让这些局部技能在共享参数里被复用，而不是为每个任务单独从零学一遍。
\end{knowledgebox}

\subsection{从 BC-Z 到 OpenVLA：任务条件的现代输入方式}

slides 还特意给了两类例子，说明 task conditioning 近年来如何与 foundation model 融合：
\begin{itemize}
    \item BC-Z 一类方法用明确的 task embedding 或多模态条件，追求 zero-shot 任务泛化。
    \item OpenVLA 这类系统则更直接，把任务描述作为 prompt 输入到视觉-语言-动作模型中。
\end{itemize}

\begin{figure}[H]
\centering
\includegraphics[width=0.9\textwidth]{slides-images/slide-19.jpg}
\caption{BC-Z：通过任务条件实现跨任务零样本泛化}
\end{figure}

\begin{figure}[H]
\centering
\includegraphics[width=0.9\textwidth]{slides-images/slide-20.jpg}
\caption{OpenVLA：把任务描述作为 prompt 输入 VLA 模型}
\end{figure}

\begin{figure}[H]
\centering
\includegraphics[width=0.9\textwidth]{slides-images/slide-21.jpg}
\caption{多任务 imitation 的机器人例子：从 ``折毛巾'' 到 ``挂毛巾''}
\end{figure}

\begin{figure}[H]
\centering
\includegraphics[width=0.9\textwidth]{slides-images/slide-22.jpg}
\caption{从少量任务到开放家务任务：多任务策略的规模正在快速上升}
\end{figure}

\subsection{本章小结}

多任务 RL 的核心不是 ``在一个网络里塞进更多任务''，而是通过条件化把不同任务写成同一策略要解决的一组相关问题。只要任务之间存在共享结构，数据复杂度就有机会被多任务训练摊薄。

